\chapter{Solutions to the Exercises}\label{ch:solutions}

Solutions are grouped by chapter. Each exercise statement links to its
solution, and each solution links back to its statement.

% =====================================================================
\section{\nameref*{ch:affine-spaces}}

\solution{affine-spaces:exr:1}
By Chasles' relation, $c - a = (b - a) + (c - b)$ and
$c - a = (d - a) + (c - d)$, so
\[
  (b - a) - (c - d) = (d - a) - (c - b) .
\]
One side vanishes if and only if the other does, which is the first claim.
For the second, the midpoints $\frac12a + \frac12c$ and $\frac12b + \frac12d$
coincide if and only if their difference is zero. This difference is the
combination $\frac12a - \frac12b + \frac12c - \frac12d$, whose weights sum
to $0$: it is a vector, and computing it with the origin $d$ gives
$\frac12\bigl((a - d) - (b - d) + (c - d)\bigr)
= \frac12\bigl((a - b) + (c - d)\bigr)$. It vanishes if and only if
$c - d = b - a$.

\solution{affine-spaces:exr:2}
Count the total weight on the points
(\cref{affine-spaces:prop:affine-comb}).
$a + 2u$: a point (point plus vector).
$\frac12(a + b)$: weights $\frac12, \frac12$, total $1$, a point (the
midpoint).
$a + b - c$: total $1$, a point (the fourth vertex of the parallelogram
with vertices $a$, $c$, $b$).
$a - 2b + c$: total $0$, a vector (twice the displacement from $b$ to the
midpoint of $[a, c]$).
$2a - b$: total $1$, a point (the reflection of $b$ through $a$).
$a + b$: total $2$, undefined.
$\frac13(a + b + c) - a$: weights $-\frac23, \frac13, \frac13$, total $0$, a
vector (from $a$ to the centroid).

\solution{affine-spaces:exr:3}
\begin{enumerate}[nosep]
  \item Fix $b\in\Bff$ and let $W = \{c - b : c\in\Bff\}$. As in the proof of
        \cref{affine-spaces:prop:subspace-comb}, $0\in W$, and if
        $w = c - b\in W$ then $b + tw = (1 - t)b + tc$ is on the line
        through $b$ and $c$, so $tw\in W$ for every $t\in\R$. Let
        $w_1, w_2\in W$. Then $2w_1, 2w_2\in W$, so $b + 2w_1$ and
        $b + 2w_2$ are in $\Bff$, and so is their midpoint. Computed with
        the origin $b$, this midpoint is
        $b + \frac12(2w_1) + \frac12(2w_2) = b + w_1 + w_2$. Hence
        $w_1 + w_2\in W$, and $W$ is a linear subspace.
  \item Fix $o$ and let $L(v) = f(o + v) - f(o)$. The identity
        $o + tv = (1 - t)o + t(o + v)$ gives $L(tv) = tL(v)$ as in the
        proof of \cref{affine-spaces:prop:map-comb}. For additivity,
        $o + u + v$ is the midpoint of $o + 2u$ and $o + 2v$, so
        \[
          f(o + u + v) = \tfrac12f(o + 2u) + \tfrac12f(o + 2v)
          = f(o) + \tfrac12L(2u) + \tfrac12L(2v) = f(o) + L(u) + L(v) ,
        \]
        where the midpoint in $\Bff$ was computed with the origin $f(o)$.
        So $L$ is linear, and $f$ is affine by
        \cref{affine-spaces:prop:one-point}.
\end{enumerate}

\solution{affine-spaces:exr:4}
Let $m = \frac12b + \frac12c$. By
\cref{affine-spaces:prop:associativity},
$\frac13a + \frac23m = \frac13a + \frac13b + \frac13c = g$. The weights
$\frac13$ and $\frac23$ are non-negative, so $g$ is on the segment
$[a, m]$, and $g - a = \frac23(m - a)$. The expression of $g$ is symmetric
in $a, b, c$, so the same holds for the other two medians: all three pass
through $g$.

\solution{affine-spaces:exr:5}
\begin{enumerate}[nosep]
  \item From $x$ and $x' = x - \omega$, the first computation gives
        $x_1 = (1 - \eta\lambda)x + \eta\gamma(x)$, whose new coordinates
        are
        \[
          x_1 - \omega = (1 - \eta\lambda)(x' + \omega) + \eta\gamma'(x')
          - \omega
          = \bigl[(1 - \eta\lambda)x' + \eta\gamma'(x')\bigr]
          - \eta\lambda\omega .
        \]
        The bracket is the second computation. The two new points differ
        by the vector $\eta\lambda w$.
  \item $x_k = (1 - \eta\lambda)^kx_0\to0$ and
        $x'_k = (1 - \eta\lambda)^kx'_0\to0$, since
        $\abs{1 - \eta\lambda} < 1$. In the original coordinates the
        second limit is the point $\omega$. Each sequence converges to its
        own origin.
  \item The map is
        $x\mapsto(1 - \eta\lambda)x + \eta(\chi - x)
        = (1 - \eta(1 + \lambda))x + \eta\chi$, affine with linear part
        $\rho\,\id$, $\rho = 1 - \eta(1 + \lambda)$. Its fixed point solves
        $\eta(1 + \lambda)x = \eta\chi$, so $x^* = \chi/(1 + \lambda)$,
        which is the barycentre $\frac1{1 + \lambda}c
        + \frac\lambda{1 + \lambda}o$ since $o$ has coordinates $0$. By
        \cref{affine-spaces:prop:fixed-point},
        $x_k - x^* = \rho^k(x_0 - x^*)$, which tends to $0$ if and only if
        $\abs\rho < 1$, that is, $0 < \eta < 2/(1 + \lambda)$ (or
        $x_0 = x^*$). With $\lambda = 1$ and $\chi = (3, 1)$,
        $x^* = (1.5, 0.5)$. If the decay is centred at the point $(2, 2)$,
        the same computation in the frame with that origin gives the
        midpoint of $(3, 1)$ and $(2, 2)$, namely $(2.5, 1.5)$.
\end{enumerate}

\solution{affine-spaces:exr:6}
\begin{enumerate}[nosep]
  \item $(f\circ\tau_\omega)(x) = W(x + \omega) + b = Wx + (b + W\omega)$:
        same weight matrix, bias $b + W\omega$.
  \item $(f\circ\tau_\omega)(x) = Wx + W\omega$ is of the form $W'x$ if and
        only if it vanishes at $x = 0$, that is, $W\omega = 0$ (and then
        $W' = W$).
  \item Passing from Celsius to kelvins replaces each input $x$ by
        $x + \omega$ with $\omega = 273.15$. With a bias, by (1) the set
        of functions of the physical temperature that the model can
        represent is the same in both units, so least squares selects the
        same function (when the minimizer is unique) and the predictions
        agree. Without bias, the model in Celsius can only represent
        functions that vanish at $0^\circ$C, and the model in kelvins only
        functions that vanish at $0$~K. By (2) these two families share
        only the zero function when $n = 1$, so the two fits differ
        unless both are zero.
\end{enumerate}

\solution{affine-spaces:exr:7}
\begin{enumerate}[nosep]
  \item If $y - x\in\Z$ then $(y + v) - (x + v)\in\Z$, so $[x + v]$ depends
        only on $[x]$. Axioms 1 and 2 hold because they hold in $\R$.
        Axiom 3 fails: $[0.1] + v = [0.3]$ for every $v\in0.2 + \Z$, so the
        displacement from $[0.1]$ to $[0.3]$ exists but is not unique.
  \item Replacing $x$ by $x + m$ and $y$ by $y + k$ with $m, k\in\Z$
        changes $(x + y)/2$ by $(m + k)/2$, which is an integer when
        $m + k$ is even and a half-integer otherwise. So the class
        $[(x + y)/2]$ takes exactly the two values $[(x + y)/2]$ and
        $[(x + y)/2 + \frac12]$: two diametrically opposite points of the
        circle. With angles measured in degrees, this is the pair
        $0^\circ$ and $180^\circ$ found for $350^\circ$ and $10^\circ$.
\end{enumerate}

\solution{affine-spaces:exr:8}
\begin{enumerate}[nosep]
  \item Solving $y = Ax + \beta$ gives $x = A^{-1}y - A^{-1}\beta$, an
        affine map with matrix
        \[
          \begin{bmatrix}A^{-1} & -A^{-1}\beta\\0 & 1\end{bmatrix} ,
        \]
        and the product of this matrix with
        $M = \begin{bsmallmatrix}A & \beta\\0 & 1\end{bsmallmatrix}$, in
        either order, is the identity:
        $A^{-1}\beta - A^{-1}\beta = 0$ and
        $-AA^{-1}\beta + \beta = 0$.
  \item $T_\omega = \begin{bsmallmatrix}I & \omega\\0 & 1\end{bsmallmatrix}$
        sends the new coordinates $x'$ of a point to its old coordinates
        $x = x' + \omega$. Then
        \[
          MT_\omega
          = \begin{bmatrix}A & A\omega + \beta\\0 & 1\end{bmatrix} ,
        \]
        the matrix of $x'\mapsto Ax' + (\beta + A\omega)$. On vectors,
        $T_\omega$ acts as the identity, since their last coordinate
        is~$0$.
\end{enumerate}

\solution{affine-spaces:exr:9}
Let $a, b\in g^{-1}(C)$ and $t\in[0, 1]$. By
\cref{affine-spaces:prop:map-comb},
$g((1 - t)a + tb) = (1 - t)g(a) + tg(b)$, which is in $C$ because $C$ is
convex. So $g^{-1}(C)$ is convex, and
\[
  f\bigl(g((1 - t)a + tb)\bigr) = f\bigl((1 - t)g(a) + tg(b)\bigr)
  \leq(1 - t)f(g(a)) + tf(g(b)) .
\]
The map $\theta\mapsto X\theta - y$ is affine from $\R^n$ to $\R^m$, with
linear part $X$, so $\theta\mapsto\varphi(X\theta - y)$ is convex. Least
squares is the case $\varphi(r) = \norm r^2$.
